\section{Certified empty-pattern handle cancellations and cocore discs}
\label{sec:source-handles}

The inherited core of Section~\ref{sec:parallelity-core} need not be
admissible. We now implement an initial geometric reduction for the empty
boundary pattern, together with a new source-bound positive unknot
certificate. The reduction removes eligible two/three-handle pairs and
one/two-handle pairs using their actual attaching occurrences. An exposed
one-handle cocore is accepted as a compressing disc only after an exact
boundary-homology check. The operation is polynomial in the source
triangulation size and does not enumerate normal surfaces.

This is one initial stage of the intended hierarchy. Its input is the
validated finite compact orientable source with one torus boundary and
empty pattern. It does not discard a pattern on a later hierarchy state,
and its output does not assert full admissibility. The implementation is
\code{reduce\_source\_handles}; the recognizer can try it with
\code{normal\_seed\_handles=True}, or the command-line combination
\code{--normal-seed --normal-seed-handles}.

\subsection{Source attaching occurrences and geometric cancellations}

The source dual structure has $T$ zero-handles, one one-handle for each
internal face pair, one two-handle for each interior edge, and one
three-handle for each interior vertex. Every two-handle carries its cyclic
edge-link sectors and its oriented passages through the one-handles.
Repeated occurrences are retained. Its contacts with three-handles are the
interior endpoints of its source edge, counted with multiplicity. In
particular, an edge looping at an interior vertex has two such contacts.

The two local moves are the empty-suture versions of Procedures 6 and 2,
respectively, in Section 6.1 of~\cite{lackenby-parallelity}. A two-handle
meeting the remaining three-handles in exactly one disc cancels with that
three-handle. After all three-handles are gone, a one-handle meeting the
remaining two-handles in exactly one disc cancels with the incident
two-handle. These are geometric incidence conditions.

For the first phase, a queue tracks source edges with exactly one remaining
interior endpoint. Removing a vertex handle exposes further edges to the
boundary. Every interior vertex is reachable from a boundary vertex in the
source one-skeleton, so this phase removes all three-handles. This is the
source-specialized realization of the published removal of three-handles
in a manifold with no closed component.

For the second phase, occurrence counts are maintained at every remaining
one-handle. When a count is one, remove that handle and its unique incident
two-handle, then decrement the counts along the removed attaching cycle.
Each source attaching occurrence is removed at most once. No remaining
two-handle uses a cancelled one-handle, so the remaining attaching cycles
retain their original source charts. Three-handles have already been
removed, avoiding a need to update their attaching data in this phase.

An algebraic coefficient is insufficient. Oppositely oriented occurrences
can cancel in a boundary matrix while remaining distinct attaching discs.
Also, a coefficient of one in a single relator does not authorize deleting
another relator's passage through the same one-handle. The checker accepts
exactly one occurrence among \emph{all} remaining two-handles; it neither
freely reduces attaching words nor supplies implicit handle slides.

The original source and the move sequence reconstruct the reduced handle
presentation geometrically. The result also lists its remaining handle
identifiers and characteristic zero-handle indices. For empty pattern, an
end of a one-handle with $d$ incident two-handle occurrences has index
$-4+2d$. At the producer's fixed point, no one-handle has $d=1$. Handles
with $d=0$ still require attention. For the empty pattern, the second
admissibility condition is automatic: any simple closed curve on a spherical
zero-handle boundary bounds a disc whose intersection with the pattern is
empty. Consequently, a replayed result in which every reported index is
nonnegative is admissible for this empty-pattern input. This observation does
not extend automatically to later states with nonempty patterns.

\subsection{Exposed cocores and the essential-boundary obligation}

A remaining one-handle disjoint from all two-handles has a properly embedded
disc cocore: its boundary lies on the manifold boundary. A graph path joining
its two ends without using that one-handle certifies that cutting along the
cocore leaves the manifold connected. The graph is the zero/one-handle
incidence graph; higher handles do not join its connected components.
A forest that processes used one-handles before unused ones supplies
candidate unused edges outside the forest and connecting paths for them.

Nonseparation alone does not prove that the disc boundary is essential in a
general reducible source. A nonseparating sphere can give a properly embedded
nonseparating disc with inessential boundary. The source operator therefore
checks the boundary cohomology class of each proposed cocore. Its positive
claim does not assume irreducibility.

\subsection{Transporting a cocore class through the cancellations}

Give the final unused one-handle value one and all other remaining
one-handles value zero. These are the values of the cocore's dual integral
cohomology class on the reduced dual graph. Reverse a one/two cancellation.
If the removed one-handle is $f$ and the removed two-handle has oriented
attaching cycle $r$, then $f$ occurs exactly once, with sign
$\epsilon\in\{1,-1\}$. Its original value is forced by
\[
 x_f=-\epsilon\sum_{(g,\sigma)\in r,\;g\ne f}\sigma x_g.
\]
This reconstructs the class on all original one-handles. Two/three
cancellations do not add one-handle values. Every original two-handle
equation is checked afterwards, including those removed with three-handles.
Thus the packet contains a source cocycle, rather than only a function on
the final graph.

The verifier uses a different characterization. It checks the prescribed
values on all remaining one-handles and all original two-handle equations.
The cancelled one-handle values are then uniquely determined by the
triangular sequence of unit equations. It does not call the producer's
reverse-substitution routine.

\subsection{A source boundary period without normal-disc expansion}

For each source tetrahedron $t$ and vertex $v$, construct an integer potential
$u_{t,v}$. Across a canonical paired face from $t$ to $s$, with permutation
$p$ and dual value $x_f$, impose
\[
 u_{s,p(v)}-u_{t,v}=x_f\qquad(v\text{ on the paired face}).
\]
The components of this corner graph are the source vertex stars. Their
links are spheres or discs, and the source edge equations give consistency.
The constructor detects any conflicting potential assignment, and the
verifier checks every face equation directly.

On an oriented source edge from local vertex $a$ to $b$ in tetrahedron $t$,
put $y_e=u_{t,b}-u_{t,a}$. Face equations make this well defined on global
oriented edges, and its sum around every source face is zero. Choosing
reference corners in the vertex stars identifies this primal cocycle,
up to the overall orientation sign, with the transported dual cocore class:
a path through a vertex star records the potential difference, and the
within-tetrahedron segment contributes no dual-face crossing.

Restrict $y$ to the boundary triangulation. A spanning-tree potential test
finds a closed boundary edge walk $\gamma$ with
\[
 \langle y,\gamma\rangle\ne0,
\]
or establishes that this candidate's restriction is exact. The certificate
records the dual values, all corner potentials, the oriented boundary walk
and its nonzero period. Global edge orientation is taken from the actual
source edge representative, not from an arbitrary final incidence.

\begin{proposition}
An accepted cocore certificate with nonzero boundary period gives a properly
embedded disc with essential boundary on the source torus.
\end{proposition}
\begin{proof}
The geometric cancellation trace identifies the final handle presentation
with the source manifold. The unused one-handle cocore is a proper disc.
Its transported dual cohomology class is determined by the unit cancellation
equations. The corner-potential equations identify its restriction to the
boundary with $y$, up to sign. Evaluating this restriction on a boundary
curve is the algebraic intersection with the disc boundary. If that boundary
were inessential on the torus, it would be null-homologous and every such
intersection would vanish. The certified nonzero period rules this out.
\end{proof}

For a source known independently to be the exterior of the input knot, this
essential disc is an unknot witness. The new
\code{diagram-handle-disc-v1} wrapper checks the original planar diagram,
reconstructs its canonical exterior independently, and verifies the entire
handle and boundary-period certificate. A reducible torus-boundary manifold
is not accepted as a knot exterior merely because it has a compressing disc.

\subsection{Size, verification and integration}

The sum of original two-handle attaching-cycle lengths is at most $6T$.
There are $O(T)$ cancellations. The queue phase uses linear incidence updates,
with source indexing and ordered traversal contributing at most
$O(T\log T)$ combinatorial work. For a candidate cocore, reverse substitution
can increase the largest absolute value by at most the length of the
removed cycle at each step. Since each such cycle is used once and their
lengths sum to $O(T)$, all transported integers have $O(T)$ bits. Corner
potentials and the boundary period obey the same bit-length bound.

One candidate uses $O(T)$ arithmetic operations on these integers.
Testing at most $O(T)$ candidate cocores therefore takes at most $O(T^3)$
bit operations with elementary integer arithmetic. A positive certificate
uses $O(T^2)$ bits. These bounds concern the new cancellation/cohomology
kernel; the established source validator and independent raw-cell
reconstruction are additional polynomial operations. The intentionally
simple independent cancellation checker rescans remaining attaching
occurrences at each proposed one/two cancellation.

Source replay reconstructs its handles from the independent raw cut-cell
model at zero normal coordinates. It does not call the source-handle
constructor, scheduler or cocore producer. It verifies exact occurrence
multiplicities, the remaining handle sets, the graph path and every
cohomology equation. Move-capped source discovery is inconclusive and emits
no partial disc certificate. A valid reduction with no disc is not a negative
recognition result. Its checked index list determines whether it is already
admissible for the empty pattern; otherwise the unused one-handles remain
explicit obligations.

The native normal-seed search can run this stage before extracting a cocycle.
On a miss it continues its existing search. Construction and positive replay
share the existing work allowance, and public certificate dispatch and the
command-line interface accept the new proof type. The stage is optional:
the existing hard survivor corpus supplies no cocore hits from these two
cancellation rules alone. General boundary-pattern transport, the remaining
admissibility reductions, and bounded hierarchy search remain obligations.

An unused bridge has a separating cocore. Its boundary is inessential on the
single boundary torus: an essential torus curve is nonseparating and would
join the two sides through the boundary complement. In an irreducible knot
exterior, such a disc cuts off a ball. This identifies the next empty-pattern
reduction to implement, together with its source-chart transport. It cannot
be applied to an arbitrary reducible source merely from the Euler
characteristic of that side. The original diagram provenance, or another
independent irreducibility justification, must remain part of the authority
for discarding it.

The source audit checks 193 triangulations: the 48 saved sources, three
relabellings of each, and an additional trefoil-exterior connected sum with
$S^2\times S^1$. Independent replay accepts every reduction. Regina's
compressing-disc oracle agrees with all 176 positive and 17 negative
source outcomes. The reducible solid-torus connected sums remain positive
for a compressing disc while not being solid tori. The reducible trefoil
control remains negative.

A further 57 canonical diagram exteriors include the hard survivor corpus,
knotted controls, the empty circle and seeded small random braids. The stage
produces 15 diagram-bound disc certificates; all agree with the independent
recognition controls, and Regina confirms the positive source exteriors are
solid tori. No hard survivor gains a cocore certificate. Thirteen of the 193
source reductions have all characteristic indices nonnegative and hence
already satisfy empty-pattern admissibility. None of the 57 canonical
diagram reductions does: unused bridge handles remain to be treated.

The full Regina-enabled suite passes all 1,626 tests with no skips. Twelve
new tests cover source reconstruction, geometric multiplicities, exact
boundary periods, forged proofs, relabellings, cancellation, shared caps,
public recognition and command-line dispatch.

The following comparison measures the complete native positive-certificate
stage from its planar diagram: exterior construction, discovery, internal
positive replay, a second public certificate replay and serialization. The
existing cocycle path and the optional handle-first path must return the
same status. Five shuffled paired rounds follow a warmup, with an identical
existing-path control. This does not time the full recognizer portfolio.

\begin{center}
\begin{tabular}{lrrr}
\toprule
Case & Existing (ms) & Handles (ms) & Existing/handles \\
\midrule
Two-crossing circle & 10.87 & 35.14 & 0.310 \\
Eight-crossing circle & 44.66 & 253.40 & 0.177 \\
Span-optimized positive & 53.66 & 120.26 & 0.443 \\
Survivor 06 & 85.27 & 110.48 & 0.791 \\
Trefoil & 15.31 & 20.83 & 0.737 \\
\bottomrule
\end{tabular}
\end{center}


The handle stage is slower on every measured case. Its three direct cocore
hits cost 2.26--5.66 times the existing stage, and the two misses add about
26--36 percent. Identical control ratios range from 0.980 to 1.008. These
measurements justify retaining the optional setting. The current independent
checker reconstructs the zero-coordinate raw-cell core and rescans remaining
attaching cycles at each cancellation; reducing that replay overhead is a
separate optimization opportunity. The present gain is a verified geometric
reduction and an additional polynomial certificate construction, not a claimed
portfolio speedup.

